% 来源：OEBPS/Text/Section0001_split_002.xhtml#sigil_toc_id_23

\BookChapter{14}
\section[第14章\q 投资营销谎言9：“我们自欺欺人的谎言”]{第14章\\ 投资营销谎言9：“我们自欺欺人的谎言”}

\begin{BookQuote}
寻找真相，你才能找到出路。

\hfill ——弗兰克·斯劳特
\end{BookQuote}

好吧，我们实话实说吧。我们前面讲了八大投资营销谎言，多年来，让我们付出了巨大的成本，而让那些大型金融机构得到了巨大的好处。我敢打赌，你现在可能非常震惊，但是你感觉内心充满了巨大的力量。你现在知道应该避开什么、做什么，才能投资成功了。

但是，你还需要粉碎最后一个投资谎言。这个谎言说的是，我们没有成功，我们没有成就，我们没有成长，都是因为别的人或者别的事情超出了我们的控制范围，我们自己没有办法解决。或者我们认为，自己天生不是那块料，根本不能掌控这一领域。但是真相是这样的：我们大多数人的人生无法实现重大进展，阻止我们的终极因素并不是别人的限制，而是我们自己限制自己的感觉或信念。不管我们现在多么成功，不管我们现在的层次有多高，个人、专业、精神、情绪层面达到了多高的层次，总还有更高的层次。我们必须对自己非常诚实，对自己没有意识到的恐惧非常诚实。我说的意思是什么？

每个人都在某种程度上害怕失败，有时我们非常害怕自己美中不足，甚至尽管我们知道该做什么，我们的恐惧害怕还是会阻止我们去执行既定的计划。结果是我们没有直接面对我们自然的恐惧害怕，我们做了什么？我们开始编故事。故事讲的是，为什么我们还没有足够聪明、足够成功、足够苗条、足够富有、足够爱别人或者让人爱。我们编的故事总是把自己的失败归因于某个超过我们控制范围之外的东西，或者我们天生缺乏某种才华和能力。但是才华和能力这两个成功的关键因素，是每一个真正承诺全身心投入的人都可以得到的。你也可以得到你想要的技能，只要你能超越那些认知上的自我限制——认为自己要掌握这些技能太辛苦、太困难，自己“根本不可能做到”。

你已经做出了一生中最重要的独立的投资理财决策，那就是精确地决定你要储蓄起来的收入的具体比例，以此来打造你的自由基金，好让你能够借助这笔资金来创造一个赚钱机器，让你睡觉也能照样赚钱。我们已经花费相当多的时间看穿那些投资营销谎言，正是由于它们的存在，我们一路走来才会不断地犯错，不断地栽跟头。那么还有什么其他障碍？最后一个阻挡我们前进的障碍是，我们自己的故事，我们自己的限制，我们自己的恐惧。我们最后要面对的那个障碍就是我们自己。正是因为这个原因，过去38年来我最有激情去做的事情就是，帮助人们冲破阻碍，帮助他们走向他们想要去的地方，而且速度更快。我的整个人生一直致力于帮助人们创造突破，而且坦白地说，尽管很多人把这一步搞得非常复杂，我却发现其实只有三个关键因素决定了你的人生是成功的还是失败的，决定了你是原地踏步还是大步向前，决定了你是否总是找借口以及是否能够享受你值得拥有的人生。

\subsection{突破}

什么是突破？突破就是从不可能变成可能的那个时刻——你不再只谈论某件事，而是终于采取大量的行动，做任何需要做的事，让这件事真正发生，这个时候就是突破。你做出一个举动，真正改变和改善你的世界，这就是突破。

经常是挫败、生气、压力引发了突破。我们到达了临界点：在这个临界点，我们会说“再也做不到，再也多不了了”。或者灵感一下子降临了：我们遇见某人激发了我们，让我们看到人生完全可以更加伟大，比我们梦想的都更加伟大。你遇见一个人，他能够充分享受人生，有很棒的人际关系和健美的身材，有真正的财务自由。这让你下定决心：“我要跟他一样聪明能干，我也能过上这样的生活。我要找到一条新的人生之路。”原来你可以接受的，现在却再也不能接受了。你现在再也不会回到从前那个样子了。你下定决心，在沙滩上画下一条线，承诺去追求一个新的目标，设定一个新的标准，从此之后，你能够做到的事会让你自己感到震惊。

大多数人会说：“这花了我十年才做出这个改变。”但是真相是并不需要花上十年，你就可以做出突破。真正的转变就发生在那一刻。也许要花费十年，你才能到达那个关键点，你已经准备好了，你放开了，甚至你被人激怒了。但是我们的人生都曾有过突破，这些突破都发生在那短短的一刻。有件事让我们挣扎了好多年——工作或者事业，体重或者人际关系，我们一直很难受、很痛苦、备受折磨，直到有一天有件事扣动了扳机。突然，“就这样了”。

“我爱你。”

“我退出。”

“我加入。”

“我们开始吧。”

不是一天，也不是一个小时，就在那短短的一刻，你的人生改变了，永远。

你是否经历过一段勉强维持的合作关系，尽管你知道自己并不快乐，你的合作伙伴也不快乐？你到了非处理这段关系不可的地步了，但是你害怕未知，害怕改变，害怕孤单，结果这些害怕阻止了你采取行动，后来你就习惯了，安逸了，再也不想改变了。

不管你让自己挣扎的是什么，我知道有一个地方，你以前有过突破。花上一点儿时间，你肯定能想出来一个地方。有哪个地方曾经一直让你很纠结——每天、每周、每月，甚至好几年，十几年，直到有一天你触碰到了临界点？你突然受到启发，或者再也忍不下去了，终于做出了一个真正的决定，要永远、彻底地改变它！你采取大规模行动，而且是马上行动来做出改变。你成功了。你最终改掉了坏习惯，你再也不抽烟了。或者你离职了，再也不干那些让你一直非常痛苦的工作了，你自己创业了。也许你最终决定开始锻炼，改变自己的体形了，或者让自己跳出那段非常糟糕的人际关系。

我希望你能够拥有那种突破。有一段时间事情看起来根本不可能改变，但是你做到了——是你让本来不可能做成的事情成功了。你绝对有能力改变你人生中的任何事情。不管这件事持续了多久，你在顷刻间决定完全彻底地改掉它，那一刻就是真正下定决心的时刻，那个决定就是你行动的依据。这就是你有过的真正突破。现在就有一个新的突破正在等着你。

\subsection{三步创造突破}

只需三步就可以创造突破：三种力量结合在一起，能够大大地改变你生活的任何一个方面。每种力量都可以单独地起作用。如果能把三种力量结合在一起，你绝对能够改变你选择的那个方面。

美国人面临的三个最大的挑战是什么？哪三个领域被事实一次又一次地证明给美国人的生活带来了痛苦？我们的投资理财，我们的夫妻关系，我们的身体健康。第一个领域是钱。有多少人，你知道，一直在钱的事上挣扎。他们不能储蓄；他们赚的钱太少；他们花的钱太多；他们想不出来为了让自己的职业发展得更好，下一步该去做点什么。第二个领域是怎么处理好自己的夫妻关系。男人和女人天生不同，如果我们不能彼此理解的话，我们就需要做很多工作，才能保持健康的亲密关系，才能理解我们的伴侣真正需要什么，真正想要什么，才能够用爱和支持的方式进行沟通。第三个领域是我们的身体。在我们现在生活的时代，西方世界的大多数人都体重超标。10个美国人里有7个人要么超重（美国疾病预防控制中心的衡量标准是身高体重比为25.0\textasciitilde{}29.9），要么肥胖（身高体重比超过30.0，甚至更高）。让我们挣扎痛苦的是体重和健康，它们已经成为全国性的危机，现在正在扩展到全世界，因为有些发展中国家也吸取了部分美国人的生活方式和饮食结构。

为什么我把那些问题全都摆出来？夫妻关系的挑战、不健康的饮食习惯和你获得财务自由的能力有什么关系？不管你想在哪个领域创造突破，不管是你的身体，你的夫妻关系，还是本书关注的焦点——钱，你都只有三样东西需要关注。你如果想改变你的人生，就必须改变你的策略，改变你的故事，改变你的状态。让我们从策略开始，因为这是大多数人开始的地方。

\subsection{正确的策略}

你一直在读本书，你肯定在寻找答案，寻找策略来掌控你的金钱，确保你未来良好的财务状况。我活着就是要找到策略来改善我们人生的每一个方面。过去我花了38年，一直执着地坚持寻找策略和工具来快速改变人们的生活质量。我做得非常成功，影响了100个国家的5000多万人。因为我疯狂地找简单的策略来快速引发突破——人生各个重要方面的突破：夫妻关系的突破，投资理财的突破，工作事业的突破，企业增长的突破，身体、心灵、灵魂的突破。

我始终相信，得到一种结果的最好和最快的方法就是找到一个榜样，这个榜样已经做成了，你只要模仿他的行为就行了。你如果知道有个人以前很胖，但是现在身材苗条，十分健康，而且保持了十几年，那么你赶紧模仿这个人的生活模式！你有个朋友，过去夫妻关系非常紧张，现在夫妻二人的感情很好，相亲相爱，而且保持十多年了，赶紧去模仿他！你碰到一些人，开始的时候他一无所有，后来一路打拼，积累了很多财富，而且财富一直很稳定。快去学习这些人的策略，这些人并不只是因为运气好，还因为他们做的事与众不同，所以他们才会在这个方面取得成功。

我全力以赴地寻找那些卓越人士的过人之处。为了找到一种行之有效的策略，你必须去找那些做得最优秀的人，他们的成功久经考验。你如果追随他们的策略，播下了同样的种子，就会收获同样的果实。我经常说“成功有迹可循”，就是这个意思。本书有很多有效的成功策略，都是从最成功的人那里总结出来的。

正确的策略的另一个好处是，可以节省你最宝贵的资源——时间。如果你一开始就有一个久经考验、行之有效的计划，有个正确的策略，那么你可以把别人几十年的挣扎奋斗，变成几天内就能实现的目标。如果你没有正确的策略，那么挫折根本无法避免，因为你第一次学某样东西，只能不断地尝试，不断地失败，不断地摸索。相反，按照正确的策略，你得到想要的结果的时间只需要几天，而不是几年，因为你学习那些已经取得成功的人，你找到了捷径。为什么要重新经历一次挫败？

现在就是关于策略力量的问题了。继续读下去，你就会拥有你在全世界能够找到的最好的投资理财策略。我可以向你保证这些策略并不是我发明出来的，而是有史以来最成功的投资者所采用的策略。但是尽管我对策略非常着迷，我也知道仅是策略本身不足以让你成功。

为什么仅靠策略是不足够的？有一个关键的挑战：人们往往采用错误的策略，不可避免地，最终会以失望而告终。你正在努力减轻体重，你用的办法是每天只吃500卡路里的食物——这肯定没办法持续，你会饿死的。或者你非常肯定自己只要买一只热门股就可以发大财了，实际上这样做几乎不可能成功。

大多数人会去找谁学习策略？我们会找谁寻求建议和指导？我们往往从失败的人身上学习！你有多少次是从夫妻关系非常糟糕的人那里寻求搞好夫妻关系的建议的？你有多少次是向一个一直为超重而挣扎的人寻求保持身材的建议的？有多少人接收的信息反而强化了他们永远无法改变自己身材的错误想法？为什么要听这种只有反作用的信息？因为没办法，他们周围都是那些身材不好的朋友或者家人。我们寻求投资理财的建议也是同样的情况。我们向那些并没有真正积累到相当多财富的人求教，因而得到的投资建议不是良药而是毒药，它们只会让自己陷入灾难。这只会加强你原来的信念，用什么办法都没用。其实，根本不是用什么方法都没用，只是因为你学习的是那些不成功的人的不成功的策略，错误的策略，而这些策略是无效的。

\begin{BookQuote}
不管策略看起来多么美，你还是要不时地看一下结果。

\hfill ——温斯顿·丘吉尔
\end{BookQuote}

\subsection{故事的力量}

回到我们面对的最大挑战：我们的夫妻关系，我们的身体，我们的财务。在每一个领域，我们都被困住了，原因有三个。第一，正如我们前面所说的，我们缺少正确的策略。我们都知道一对夫妇，不是男的不愿意沟通，就是女的不能停止唠叨。夫妻之间不明白对方的需要，更不用说满足对方的需要了。如果有个朋友总是不断地追逐最流行的减肥食谱，你会怎么跟她说？或者有个朋友总是在寻找一个神奇的方法能使自己迅速地成为百万财富，他其实是在告诉自己，若没有这个神奇的方法，他就永远不能实现财务自由。没有正确的策略，你就会失败。当你失败了，你就会编出一个很离谱的故事：“我的老婆永远不会满足。”“我永远不能减掉体重。”“能赚到大钱的人是那些已经很有钱的人。”这些自我选择的故事阻止我们去寻找正确的策略，甚至尽管我们有了正确的策略，它也会阻止我们具体实施。

你是不是也认识这样的朋友？你把答案直接放在他们的眼皮底下，他们还是会说：“不行，这根本不会有效，因为……”他们会告诉你100万个理由。正确的策略就摆在面前，为什么人们就是不用？为什么他们还是没有实现他们的目标？要保持相亲相爱的夫妻关系怎么就这么难？为什么彻底减掉体重这么困难？为什么70\%的美国人的体重都超标？是因为变苗条和健康的策略实际上太复杂了吗？是不是那些信息都被隐藏起来了，只有1\%的人才能找到，或者这些信息收费非常高，只有特别有钱的人才能买到？见鬼去吧，根本不是。你要找的答案随处可见，随时可得，随地可用。你家附近肯定有一个健身房，开车一会儿就到。（上帝禁止我走过去。）全世界到处都有教练，有些能在网上指导你，你不管在哪里都行！网络上到处都有免费的建议，当然，网上还有好多关于健身和减肥的书籍，你用手机或平板电脑可以马上下载。你根本不用费事就能找到你想要的策略，让自己的形体变得更好，体格变得更强壮，身体变得更健康。

那么真正的问题是什么？答案是我们必须把人性考虑进来。我总说，人生成功，八成靠心理，两成靠方法。你怎么解释有些人明明知道自己需要做什么，自己想做什么，而且有正确的策略，他却没有采取行动？要解开这个谜题，我们必须深入地研究个人的心理：价值观、信仰、情绪。这些心理因素在驱动着我们。

正确的策略就放在面前，但是你还没有成功。这是因为你丢失了达成突破的第二把钥匙：故事的力量。如果你没有采取行动，那么答案就放在你眼前，只有一个原因：你创造出了一套信念，已经绑定在一个故事里了——这个故事讲的是为什么这个方法你用不行，别人用就行。这个方法只适用于四种人：有钱的人，有好身材的人，幸运的人，夫妻关系和谐快乐的人。你很容易编出一个自我限制的故事。

既然如此，为什么你还要费事采取行动，实践你“知道”肯定会失败的策略？这个策略不是问题，故事才是问题。那些做事半信半疑的人会说：这可能能行，也可能不行……”当然会不行了！这个信念成了一个自我实现的预言。这种故事就像能够解除力量的魔戒，失败简直板上钉钉。当然，失败的结果只会进一步加强你的那个信念：用什么方法都不行。这样就形成了恶性循环——不相信，失败，更不相信，更失败。

但是那些真正做成事情的人、完成目标的人、转变的人、成长的人、学习进步的人，他们不仅采取了策略，还附加了一个新的故事：一个能够赋予人力量的故事，一个“我能做到而且我愿意去做”的故事来代替“我不能做到，所以我不愿意做”的故事。于是，实现了从一个自我限制的故事转到了一个赋予力量的故事。我不愿意成为大多数不能做到的人中的一员，我要成为少数能够做到的人中的一员。

有段时间，我的体重超标38磅，我给自己编的故事是“我的骨架很大”。确实如此，但我也很胖。故事可能是真的，但如果它们不能帮我们，而且会阻止我们拥有自己想要而且值得拥有的生活，我们就要改变这些故事。我们的人生中也有过很烂的故事。

我不能赚到足够多的钱。

我不能储蓄更多。

我从来不会阅读。我天生有阅读障碍。

我的朋友理查德·布兰森是维珍企业帝国的掌门人，他有阅读障碍，但这一点并没有限制他的人生发展。为什么？因为他对阅读障碍的信念或者故事是积极的，而不是消极的。他的故事不是“我根本不会阅读了”，而是“我有阅读障碍，所以我必须更加刻苦努力地工作，才能做成我想做的每一件事——我会更加刻苦努力”。要么你利用你的故事，要么你的故事利用你。找到你人生出错的地方，和找到你的人生做对的地方一样容易，只要你改变你的故事。你的夫妻关系出了问题，理由是所有的好男人都走光了，他们是同性恋可是你不是，或者你是同性恋可他们不是。总是有个故事好讲，对吧？故事控制着我们的情绪，情绪驱动着我们的行为和行动。

我来问你一个问题：你担心钱的事吗？钱的事会让你半夜睡不着觉，翻来覆去吗？一提钱的事你就非常紧张：下个月工资啥时候发？下个月的车贷怎么还？小孩子上大学的学费啥时候能攒够？我能不能攒到足够多的钱享受退休生活？你的财务压力到底有多大？按照美国注册会计师协会的研究，44\%的美国人，也就是接近一半的美国人，在接受调查时说自己的财务压力属于“高水平”。你有没有这样想过，“这么大的压力会压垮我”？

凯利·麦戈尼格尔是斯坦福大学的健康心理学家，他一直警告人们注意压力的危险，已经有整整10年了。后来她才意识到也许是她的建议，而不是压力本身正在把人更快地送进坟墓。“我是在把一种可以强化健康的刺激转化成一个疾病的源泉。”思路的突破加上一些强有力的新研究，麦戈尼格尔的观点完全转变了。

研究结果表明，压力也可以成为我们的朋友。就像你给肌肉压力（举重或跑步）能让肌肉更强壮，情绪压力能让我们的身体和心灵更加强壮。麦戈尼格尔现在重点普及的观点是，新的研究表明，转变你的头脑对压力的看法就能真正改变你身体对压力的生理反应。在一个持续8年的研究中，有些成年人体验到了“压力非常大”，内心也相信压力有害健康，他们的死亡风险增加了43\%。（这真让我紧张。）可是另一些成年人也同样体验到“压力非常大”，但是内心不认为压力有害健康，没有把压力看作有害的，结果他们死亡的风险一点儿也没有增加。麦戈尼格尔说，压力的生理信号（心跳加快，呼吸急促，浑身冒汗）未必是焦虑的生理证据，也未必表示我们应对压力做得不太好。相反，我们可以把这些生理反应换一种方式解读：这是我们的身体正在释放活力，让我们准备好迎接下一个挑战。现在的科学证明，关键不是压力，而是你怎么看待压力——你联想到的压力的故事很关键。告诉自己压力有害，可能会让你50多岁就得心脏病；告诉自己压力有益，可能会让你健健康康地活到90多岁。这仅是一念之差。

\begin{BookQuote}
成功是我唯一的选择，失败不是。

\hfill ——《失去自我》（歌曲）
\end{BookQuote}

那么关于钱，你过去一直给自己讲的故事是什么？是什么阻止你不能实现你的财务梦想？你一直告诉自己现在开始储蓄为时太早，还是说现在开始重建你的投资组合为时已晚呢？你现在赚的钱太少了，根本不够花，哪里还能有钱储蓄？整个社会有人操控，故意和你做对，所以你何必费事瞎折腾？也许你的故事是，“政府让我们承受沉重的债务负担，整个金融体系混作一团”，或者是“我就是对数字不在行”。好消息是，你并不需要对数字在行！你只要有个手机，上面有计算器，就可以在手机上下载一个应用程序，回答几个简单的问题——你现在在哪里，你想去哪里，你愿意做什么，从而得到一个投资理财规划。你会清楚地了解到如何实现财务自由。

也许你的故事是“先要有钱才能赚钱”。我把本书的初稿分享给一些朋友看，第一批读者中的一位，她有个核心信念是：“我永远不会实现财务自由，除非我能有办法赚到很多钱。只有一开始就有很多钱的人，才能赚到更多的钱，但是我没有那么多钱。”后来她读到西奥多·约翰逊的打造你自己的赚钱机器的故事之后——那个送快递的家伙年收入从来没有超过1.4万美元，却利用自己打造的赚钱机器，经过一生的持续努力，把这么少的收入转化成7000万美元的巨大财富，她马上把自己的故事扔到了窗外。西奥多并不幸运，他能积累如此巨大的财富，用的只是一个简单的系统，和本书后面你要学习的系统一模一样。

她的新故事也可以成为你的新故事：“我只要利用那个简单的复利机制，就能赚到很多钱。我就可以去任何我想去的地方，我就可以过我真正想过的生活，我就可以实现财务自由了。根本没有什么限制，那些限制都是我自己给自己加上去的。

我的重大财务突破之一，就是在我的故事有了重大改变后才出现的。小时候，我家很贫穷，我总是把缺钱和家里每个人的痛苦联系在一起。我很小就发誓，我绝对不会要小孩，除非我真的有钱了。我发誓，我要赚很多钱，让我的家人再也不会为钱发愁。我小的时候，因为没钱付不起账单，吃饭时桌子上什么食物也没有，我当时感到羞愧、挫败和痛苦。

我不断地努力，终于实现了自己的诺言。我18岁的时候，每个月能挣到1万美元，当时来看，这可是一大笔钱，现在也是。我非常激动，我回去找我从小一起玩到大的好朋友，他和我一样从小家里就很穷。我说：“我们一起出去大玩一场。我们飞到埃及，骑骆驼比赛，从这个金字塔骑到那个金字塔！”我还是小孩的时候就一直有这个梦想。现在我可以跟好友一块儿实现这个儿时的梦想了。可是好友的反应完全不是我想的那样，他说：“这事对你很容易，你是有钱人。”我把这家伙当好朋友，才会来找他一起去玩，没想到他竟然会这么鄙视我，我简直气得浑身发抖。我不是来炫耀我有钱的，我只是想和朋友分享我的财富，去共同体验一场真正的冒险。但是我需要重新评估，我创造出了一个新的故事：虽然一个信念说你能做好，但你做得差不多好就行了，不然别的人要说你了。如果你做得太好了，大家就会不喜欢你了。

就这样过了很多年，我的人生和我的事业都发展得挺好，但是我的收入并没有大幅增加，直到后来我最终碰到了一个临界点，到了我人生的一个重要阶段。我开始想：“这简直太可笑了。”如果我能拓展我的智力，我应该去做吗？我的答案是：“当然了！”如果我可以体验到更多的爱，给予更多的爱，我应该去做吗？当然了！如果我可以拓展我的能力，给予别人更多，我应该去做吗？当然了！如果我可以赚更多钱，增加我的财富，我应该去做吗？第一次，我感到犹豫不决。为什么在我人生的其他领域我拓展自己，让自己得到更多，给予更多，是很自然的事，但是一旦涉及钱的问题，突然之间我的答案就完全不同了？为什么？这太不合理了。

但是后来我知道真相了。我有一个深深的恐惧，自己原来并没有意识到，那就是我害怕别人指责我，于是我在这个领域进一步拓展了自己。我太希望取悦每个人了，我太希望得到别人的爱了，以至于我潜意识里不但以为赚更多钱、发更大的财像是做错事，而且它在潜意识里暗中破坏自己的成功，好让自己别赚太多的钱。像很多人一样，我告诉自己钱不是精神上的，钱和心灵无关。这是不是太疯狂了？任何一个变得真正富有的人都知道一个真相，要变得富有并且保持富有的唯一方法就是，在一个大家真正看重的领域，找到一个方法，为别人做更多，比其他所有人做的都要多。如果你成了别人生命中的“好福气”，你自己也会得到上天赐予的“好福气”。钱，只是这些“好福气”中的一种，但是它确实是一种“好福气”。它只是另一种形式的自由和富足。

钱不是别的东西，钱就是一种反映。它能反映出你的创造力有多大，你的专注能力有多强，你增加价值并且得到回报的能力有多大。你如果能够找到一种方式，创造出更多价值，让很多的人受益，那么你就有机会得到很多的经济回报，让自己的人生变得非常富足。

我到达了那个临界点，那时我开始厌烦用原来的那种方式生活，我看到了只是为了讨好别人是多么的荒谬可笑。如果你赚大钱了，别人就会把你归为“1\%的富豪精英”，这确实不假。我从还是个孩子的时候就觉得，成为那1\%精英中的一员，太野心勃勃了。我的父母属于那99\%的平凡人，我原来也属于那99\%的平凡人，但我就是不甘心这样过一辈子，为了我自己，为了我的家人，我想出人头地。不过，停留在那个层次，只是为了和别人合得来，搞好关系，这也太不合情理了。我过去把自己在财富积累上缺乏进展归咎于要和别人搞好关系，要合群，我已经厌烦透了，我再也不要这样做了。我在赚钱积累财富上限制自己的故事，必须得抛弃了。我爱别人，但是我不能浪费自己的生命来取悦他们，特别是我知道这样取悦他们，我的人生格局就小多了，我只能限制自己。我打心眼儿里不相信造物主创造我们的生命就是为了这个目的。现在是时候了，我应该去找一种方法，赚更多的钱，并用同样的方法，努力给予更多，贡献更多，付出更多的爱，同时拓展我的能力，在智力上、情绪上、精神上都要拓展自己，使自己更加精进。

有了这种信念上的转变，突然之间一切变得一清二楚了。这并不是说我要征服这个领域，比如在处理夫妻关系这个方面，策略可能一直就在那里，但是由于我的思维僵化，所以我一直对它视而不见。一旦你改变了你的故事，你的整个世界都会随之改变。

改变你的故事，就能改变你的人生；脱离那种自我限制的故事，拥抱发现真相的故事，每件事情都会改变。我敢说，当你摆脱了自我限制的故事，采取大胆的行动，你就会发现策略开始起作用了，你能创造的结果简直就像奇迹一般。

我给你最后再举一个例子，我有个好朋友朱莉，她是一个成功的编剧，她的稿酬属于顶级水平，可是她在财务上一直没有太大的进步。等到她和她的丈夫都50多岁了，他们才有了一套漂亮的房子，抵押贷款虽不算多，但是他们的个人退休金账户里只有10万美元——离他们退休需要攒的退休养老金的目标金额还差得很远。他们的钱都被投资到一个名为“社会责任”的公募基金上了，这个基金收取的费率很高，吃掉了他们的大部分投资收益。

朱莉的丈夫柯林想投资得更加激进一些，但是朱莉甚至根本不愿意和他谈投资理财的事情。朱莉告诉丈夫说他恨华尔街，恨华尔街代表的一切。事实上，一谈钱就会让她很不舒服。对于她来说，钱是罪恶的。

但是后来一个突破发生了，她参加了我举办的研讨会，释放了内在的力量，我们运用策略故事的力量改变了她的心灵状态，告诉她，你在人生的每一个方面都能创造出突破，释放出你的内在潜力。我运用音乐和其他工具来让观众处在巅峰状态，就是这个时候，正是发生突破的时候。

那个周末转变了她的财务人生。她是怎么做到的？她认识到有些事情必须改变，要不然，她和丈夫就要面临非常悲惨的退休生活，最打动她的是，她想明白了自己的负面信念，她对钱的负面看法创造出了持续不断的痛苦，由此影响了她的婚姻，也影响了她的将来。她问自己这个故事是从哪里来的，她剖析自己的内心，然后问自己，这真的是我吗？这就是我非常相信的东西吗？我们并不是生来就相信钱的，很多人认为钱是罪恶的，这个信念是从哪里来的？

朱莉不需要剖析得很深，就能找到答案。朱莉的父母成长在大萧条期间。她的妈妈从来没有得到过上大学的机会，尽管她的学习成绩非常拔尖，但她只能去找一份百货商店店员的工作，一个星期才挣9美元，而且根本不敢抱怨工资低，也不敢抱怨站得时间太长，脚疼得要命，因为她的妈妈太怕丢掉这份工作了。朱莉从小到大一直在听妈妈唠叨同样的故事，一遍又一遍——那些富人如何剥削穷人；银行和华尔街的经纪人如何摧毁了美国经济；人们根本不能相信股票市场。听的次数多了，朱莉在脑子里形成了一种联系：“如果我成了一个富有的投资者，我就变成了一个坏人，我的妈妈就再也不爱我了。”

朱莉意识到她一直在告诉自己这个讲述财富罪恶的故事，但这并不是她本人的故事，而是她妈妈的故事。“钱是万恶之源”，这是她妈妈像念经一样一天到晚重复唠叨的咒语，但并不是她自己的咒语。自己意识到了这个真相，这让她感到震惊，这个真相让她解脱了，这些话语对她失去了所有的影响力。事实上，她自己做了一番功课，研究了《圣经》上的内容，结果她发现《圣经》并没有说“钱是万恶之源”，而是说“爱钱如命”，爱钱超过爱其他所有的东西——真爱、夫妻关系、奉献，绝对会让你大难临头。

这是个令人震惊的转变。朱莉从她自我限制的故事里面走出来后，很快便坐下来跟她的丈夫第一次谈论起他们的财务问题，她非常激动，夫妻俩可以合作，重新获得财务人生的控制权。想象一下，打造你的财富会是多么艰难，因为你的核心理念是钱是罪恶的东西。他们抛弃了成本很高的公募基金，而把他们的个人退休金账户转换成一个多元化的由先锋公司提供的指数基金组合，然后他们制订了一个合适的长期投资理财规划，就像你在本书后面会读到的那样。这让他们成功地踏上了通向财务自由的阳光大道。

朱莉和柯林转变了他们的故事，结果后来发生了什么？他们学会了如何玩这个投资游戏并且获胜，他们学到了如何创造一份终身收入，他们学到了额外多加15万\textasciitilde{}25万美元到自己的腰包里的方法，它其实很简单，就是从那些收费很高的公募基金里跳出来，转向那些收费很低的指数基金，这样你的投资就能节省出很多费用。有了足够的退休金，你的退休生活会多么的幸福惬意。

记住，秘密很简单：改变你的故事，就能改变你的人生。脱离那些让你自我限制的故事，拥抱真相。你能让任何事都梦想成真。

\subsection{你的状态很重要}

如果你感觉糟透了，你根本不会认为生活很美好。你是否有过这样的经历：对某人感到很生气，然后就突然想起这个人曾经做过的每一件让你讨厌、愤怒的事情？当你进入愤怒的状态时，大脑会支持这种状态，那么让你保持这种状态的那些故事马上就会出现了。

相反，沐浴在爱情里的女孩子看世界会是什么样子？好像一切都透过了玫瑰色的透镜，每件东西看起来都美妙极了，对吗？粗鲁的店员无法惹怒你，哭闹的小孩看起来很可爱。这就是一个积极的状态，它会改变你的外表和你的故事。

你的内心和情绪状态会影响你对生活中每件事的感觉和体验。我和任何人合作，从世界一流的运动员到大权在握的公司高管，我首先要做的事就是改变他的状态。你的好心情如果被开启，你就能做成你想做的事情，但是如果被关掉了，你的整个世界就完了。你知道我说的是什么意思，对吧？你知道，当你状态很好的时候，每一件事情都进行得很完美，你甚至根本不用思考：打网球，你会发球得分；开会，你会说到点子上；谈判，你会得到你想要的结果。我们也都体验过相反的情形：我们连自己家的地址都想不起来；到别人家吃晚饭，可连女主人的名字都想不起来了；连英文单词“the”都能拼错。我称之为愚蠢的状态。但是过了几分钟你就恢复正常了：你想起答案了，因为你现在处于另一个不同的状态。

本书的目的并不是要教你如何改变你的状态。我的书、音频、课程、现场活动，很多讲的都是这个主题。总之，你可以马上彻底地改变你的感觉（而不只是希望你自己能够感觉良好），你只需要明白先改变你的身体，然后就能改变你的头脑。

我可以教给你很多办法，马上改变自己的状态。其中有一个最简单的办法，我称之为改变你的生理机能。你可以改变你思考的方式，方法是改变你的运动和呼吸的方式。情绪是动作创造的。大量的行动可能会治愈你所有的恐惧。好好想想，恐惧是生理性的。你的嘴、你的身体、你的胃都能感觉到恐惧。勇气也是一样的。你可以只用1/1000秒就从一个状态转到另一个状态，只要你能学会彻底地转变你的行动、呼吸、说话方式，以及你使用身体的方式。我运用这些诀窍近40年了，并把它们用来转变一些精英者的状态，包括世界上顶尖的运动员、交易员、企业领导人、政治领导人。2013年，哈佛大学做了一项科学研究，研究证明了这种方法的有效性。

哈佛大学教授、社会生理学家艾米·库迪在2012年发表的TED演讲非常有名，他提出了一个“非科技生活妙招”，他让在场的观众改变姿势，并持续保持2分钟。艾米·库迪的研究表明，你只要采取“有权有势的姿势”（比如电视剧《神奇女侠》中那个女超人双手叉腰，两腿坚定地站在地面上；或者想想你办公室里的那个家伙，向后靠在椅背上，两手交叉放在头后，胳膊肘向外张开——你知道那个家伙是谁），就会让睾酮（显性激素）分泌增加20\%。与此同时，皮质醇（最主要的压力荷尔蒙）分泌会减少25\%。这种荷尔蒙水平的改变影响很大，会马上改变你直面恐惧、接受风险的意愿，而所有这些变化在你的身体姿势改变2分钟后就会出现。按照艾米·库迪的研究，在摆出有权有势姿势的人中，86\%的人反馈更愿意尝试抓住机会，接受挑战。而另一组志愿者要求用更加被动的姿势站着或者坐着2分钟，双腿和双臂紧紧交叉，他们的睾酮分泌水平下降了10\%，而压力荷尔蒙水平上升了15\%。这一组志愿者的行为举止更加自信的比例要少得多，只有60\%。记住，这些不仅是生理上的变化，而且是真正的生物化学上的变化。我传授这些方法已经有38年了，我的学生们都知道我教他们的方法真的很有效，现在也被科学实验证实了。这意味着什么？大体可以说，你激动了。你的步伐迈得更大了，你准备冒险了，你要接受必要的风险来重新打造自己的世界。只需要摆上2分钟的姿势，你就可以导致荷尔蒙分泌的变化。这种变化要么让你的大脑更加肯定、更加积极、更加自信、更加舒服，要么让大脑更加被动、更加顺从。我们的身体能够改变我们的大脑。

我的人生也曾低迷过，有段时间我体重超标，心情沮丧，住在加利福尼亚州威尼斯只有一个房间的小公寓里。我盯着空荡荡的房间，听着尼尔·戴蒙德的唱片。听起来很吓人，是吧？有个朋友，我们好久没见过面了，他顺道来看我。一进房间看到我这个样子，就问：“老兄，你出什么事了？”真是一语惊醒梦中人。当时我就下定决心，我要打破现状。

于是我穿上跑鞋，抓起我的索尼随身听（不错，大叔我年纪够大才会有这个老古董）。那个时候，我沉浸在音乐里：当时，一次只能听一张唱片，不像现在用iPod（苹果公司音乐播放器），有10000首歌曲供你随意选择。我找到传奇摇滚乐队Heart（心）的唱片，播放歌曲\textbf{Barracuda}（《梭鱼》），让节奏点燃我的激情。我开始奔跑，我下定决心，要跑得很猛很快，像我过去的人生那样，不跑到吐血决不停步。若要说我下定决心要超越自己的限制，那就显得太低调、太保守了。

我敢肯定，那时我跑步的样子看起来一定非常滑稽，因为我体重超标38磅，一跑起来我的啤酒肚就被甩来甩去，活像个猪八戒。后来我累得连口气都喘不上来了，一下子瘫倒在沙滩上，手里抓着一本杂志，那是我随身带来的。就是在那种极端的信念、决心、兴奋、筋疲力尽交织在一起的状态下，我坐下来，写下了每一件我在生活中再也不会忍受下去的事情：我的身体状态，我的懒惰，我肤浅的亲密关系，我糟糕的财务状况。我决心要创造出人生新状态——在那个热烈期望、充满活力的状态，我非常肯定自己能够找到成功实现目标的道路。

有了足够强大的状态，你会发展出一个强大的故事。我的故事是：“过去的生活在此时此地结束了，新的人生从此开始了。”我非常认真，我要竭尽全力地去做。我发现，你如果改变了你的状态和你的故事，就能找到或者创造出正确的策略，得到你全心全意想要得到的东西。正是用这种方式，你能创造出一个真正的突破——一个新的状态，一个新的故事，一个得到证实的有效的策略。

过了30天，我减掉了30磅。过了6个多星期，我一共减掉了38磅。我非常狂热地努力实现我的承诺。那一天我设立了新的标准，我是什么人，我会容忍什么。从那天到现在，已经过去30多年了，这个标准一直没有被降低过（我的体重也从来没有再回到那时的水平）。

我从年收入低于3.8万美元到现在年收入超过100万美元，只用了一年多的时间。收入水平能提高这么多，我当时做梦也想不到。更重要的是我的情绪和心理状态恢复到了正常的水平——这两种力量能够真正改变某些人的人生轨迹。决心、信念、勇气开始成为指引我奋力向前的动力。

伟大的策略就近在咫尺，但你视而不见，直到你让自己进入一个合适的状态：强大有力、下定决心、自立自强。这个状态将会自动孕育出信念和故事：你能实现愿望，你必须实现愿望，你会实现愿望，而且你会全身心地投入去实现愿望。有了这样的状态和故事，你不仅能够找到有效的策略，还能拥有你值得拥有的回报。你是不是看得非常投入？只要你现在过的生活有任何一个方面远远低于你值得拥有的生活，那么现在是时间去改变这些要素中的一个或者多个了。

记住，我们能得到的东西，就是我们能容忍的东西。所以，再也不要容忍你自己给自己找的那些借口，限制自己的信念，得过且过和恐惧害怕的状态。运用你的身体，以此为工具，让你快速进入一个只有纯粹的意愿、决心和承诺的地方。直面挑战，你内心坚定的核心信念是：问题只是通往梦想道路上的减速带。从那个地方开始，你只要采取大规模的行动，按照一个有效的和经过证实的策略采取行动，就会重新改写你的故事。

现在是时候了，我再也不要成为大多数人中的一员，而要成为极少数人中的一员。成为极少数人中的一员，一步一步地提高，一步一步地进步，拥有属于自己的真正能力，不仅在投资理财上，而且在你人生中的每一个方面。大多数人一开始有很大的抱负，但是后来却勉强地接受远远低于自己真正能力的生活和生活方式。他们让失望摧毁了自己。失望是不可避免的，只要你想努力做一件大事。但是你不能投降，要让你的失望来驱动你寻找新的答案，要用纪律约束你的失望。从每一次失败中学习，按照自己从失败中学习到的经验教训去改善自己的行动。你只要这样持续做下去，成功就会是必然的。

因此，下一次当你又编出一个理由解释为什么你不能做成某件事时，而你心里知道你的精神是没有限制的，你就要直接对自己说这些理由都是胡说八道的。改变你的状态，改变你关注的焦点，重新回到真相，调整你的路子，追逐你真正想要的目标。

好了，深呼吸，或者大叫，起来，运动。把这九大投资营销谎言，这些过去的限制，全部清除。然后，我们就可以迈出下一步，也就是7步通向财务自由之路的第三步。我们要让这个游戏变成我们可以赢的游戏，其办法是想出一些具体的财富数字——这些财富数字反映出你真正想要实现的财务梦想究竟有多大。接下来，我们可以创造出一个投资理财规划，不断地改进和完善这个规划，找到办法加速实施这个投资理财规划，从而让你能够更早地实现你的财富梦想——比你能想象到的还要早。
